% Standalone uncompiled research note. No external TeX inputs are required.
\documentclass[11pt]{article}
\usepackage[T1]{fontenc}
\usepackage[utf8]{inputenc}
\usepackage[margin=1in]{geometry}
\usepackage{amsmath,amssymb,amsthm,mathtools,enumitem}
\usepackage[hidelinks]{hyperref}
\newtheorem{theorem}{Theorem}[section]
\newtheorem{proposition}[theorem]{Proposition}
\newtheorem{corollary}[theorem]{Corollary}
\newtheorem{lemma}[theorem]{Lemma}
\theoremstyle{definition}\newtheorem{definition}[theorem]{Definition}
\newtheorem{example}[theorem]{Example}
\theoremstyle{remark}\newtheorem{remark}[theorem]{Remark}
\newcommand{\Pow}{\mathcal P}
\newcommand{\Gen}{\operatorname{Gen}}
\newcommand{\Meet}{\operatorname{Meet}}
\newcommand{\cl}{\operatorname{cl}}
\newcommand{\im}{\operatorname{im}}
\newcommand{\Set}{\mathbf{Set}}
\newcommand{\Cat}{\mathbf{Cat}}
\newcommand{\Bool}{\mathbf{2}}
\newcommand{\N}{\mathbb N}
\newcommand{\NN}{\mathbb N}
\newcommand{\Cost}{\mathbb N\cup\{\infty\}}
\newcommand{\Lan}{\operatorname{Lan}}
\newcommand{\Fact}{\operatorname{Fact}}
\newcommand{\id}{\mathrm{id}}
\title{Changing experiments and composing characterizations}
\author{A research note with exact criteria and native witnesses}
\date{5 November 2026}
\begin{document}
\maketitle
\begin{abstract}
Changing the experiments or operations available to a system can alter its
behavioral quotient, concept completion, and incidence topology in different
directions and at different times. We derive these changes from native
continuation tests, prove the least simultaneous operational repair, and
characterize when repair commutes with an actual wiring. A two-bit gate
example leaves the state quotient fixed while forcing a larger concept
lattice; a subsequent native reset changes only the witness topology.
For experiment formation by declared decompositions, the comparison between
changing interfaces and composing is a pointwise Kan-extension mate.
Its established exact-square criterion leads to a complete classification
for integer resource translations: Boolean feasibility, set-valued exactness,
and equivalence of retained allocation categories require successively
stronger matrix conditions. A response-preserving cone states the precise
bridge from these experiment comparisons to the heterogeneous invariants.
Finally, synchronized weighted native machines show that complete Boolean
agreement can hide an unbounded resource defect; a common endpoint-preserving
primitive normalization gives an effective witness compiler for arbitrary
finite runs. General categorical machinery and previous project results are
credited explicitly. The concrete classifications and their integration are
proved here without a claim of historical priority.
\end{abstract}

\section{The question and the level of the answer}\label{sec:question}
An experimental characterization is relative to what can be done and what
can be observed. Suppose an old repertoire is enlarged by a gate, an
instrument, a continuation, or a changed interface. Does the old
characterization still support the new operations? If components are then
wired, does changing their characterizations first give the same result as
changing the complete system? The answer depends on both the composition
and the information retained in the characterization.

The first part of this note uses actual deterministic programs and their
Boolean responses. It derives an enlarged test doctrine, rather than
postulating a refinement map, and compares predictive states, formal
concepts and supported topology. The second part studies experiment
formation by a supplied category of decompositions. Here the relevant
Beck--Chevalley comparison can fail through an absent lift, disconnected
lifts, or a noninvertible witness that an extensional colimit forgets.
The third part makes native implementation explicit in finite weighted
machines. It separates independently optimal local implementations from
one jointly executable implementation and supplies a finite primitive
certificate for their agreement.

These are different presentations of a common problem, not interchangeable
names for one construction. In particular, FCA completion is not asserted
to be Day convolution, and CNOT is not assumed to be a Kan-extension
composition. The bridge is an evaluation map from generated experiment
labels to their actual response columns; its hypotheses are proved or
stated where it is used.

\paragraph{Reading route.}
Sections~\ref{sec:source}--\ref{sec:cnot} derive the mixed changes and a
fully computed interaction. Sections~\ref{sec:decomposition}--\ref{sec:slice}
give the categorical comparison, resource classification and response
bridge. Section~\ref{sec:native} supplies native witness conditions.
Section~\ref{sec:contextads} gives the exact contextad connection, and
Sections~\ref{sec:prior}--\ref{sec:limits} delimit the increment and the
remaining question. The source is deliberately self-contained and
uncompiled. Finite checks and intermediate research files accompany it.
